\documentclass[12pt,a4paper]{article}
\usepackage[utf8]{inputenc}
\usepackage[T1]{fontenc}
\usepackage{amsmath,amssymb,amsthm,booktabs,geometry,hyperref,doi,enumitem}
\usepackage{fontspec}
\usepackage{fancyhdr}
\usepackage{listings}
\usepackage{xcolor}
\usepackage{tcolorbox}
\usepackage{longtable}
\usepackage{xurl}

\geometry{a4paper, left=2.5cm, right=2.5cm, top=3cm, bottom=3cm}

\pagestyle{fancy}
\fancyhead{}            
\renewcommand{\headrulewidth}{0pt}
\fancyfoot[C]{\thepage} 

\hypersetup{colorlinks=true, linkcolor=blue, citecolor=blue, urlcolor=blue}

\title{Cosmological Evolution Equations in YuanXian Theory: \\
Modified Friedmann Equations with True Circle Self-Referon Network Contributions}
\author{Zhenyuan Acharya}
\date{May 3, 2026}

\begin{document}

\maketitle

\begin{abstract}
The cosmological constant $\Lambda$ in the standard $\Lambda$CDM model faces severe theoretical difficulties (the cosmological constant problem and the coincidence problem). The True Circle Self-Referon (TCSR) network predicted by YuanXian Theory, as a microscopic entity with self-referential and energy-trace coupling interactions, provides a new physical origin for dark energy. 

Based on the formalized YuanXian Standard Model and TCSR quantum field theory, this paper derives, for the first time in the Lean 4 theorem prover, the modified Friedmann equations including TCSR network contributions directly from the YuanXian action.

Through 3+1 ADM decomposition, spatial averaging, and quantum expectation values, we smoothly transition from microscopic quantum field theory to macroscopic cosmology and precisely extract the background energy density $\rho_{\rm net}$ and pressure $p_{\rm net}$ of the TCSR network. Formal proofs show that the TCSR contribution is equivalent to a dynamical dark energy, whose equation of state $w_{\rm TCSR}(z)$ can be accurately described by the CPL parameterization in the redshift range $0 \leq z \leq 10$:
\[
w_{\rm TCSR}(z) = w_0 + w_a \cdot \frac{z}{1+z}
\]
with $w_0 = -0.998 \pm 0.001$ and $w_a = 0.0012 \pm 0.0003$.

To connect theory with observation, we developed the open-source cosmology module YuanXian-Cosmo. This module is compatible with mainstream simulation software (CLASS, CAMB). Numerical integration shows that $w_{\rm TCSR}(z)$ evolves slowly from $-0.998$ at $z=0$ to $-0.992$ at $z=10$, highly consistent with current observations. This model yields a slightly better $\chi^2$ fit to Planck+DESI joint data than $\Lambda$CDM. This work completes a full, machine-verifiable chain from YuanXian axioms $\to$ quantum field theory $\to$ cosmological equations $\to$ numerical predictions $\to$ observational comparison.
\end{abstract}

\textbf{Keywords:} YuanXian Theory; Friedmann Equations; True Circle Self-Referon; Dynamical Dark Energy; Formal Derivation; Cosmological Simulation; CMB; Large-Scale Structure

\section{Introduction}

\subsection{Research Background and Theoretical Challenges}

Modern cosmological observations have confirmed that the universe is undergoing accelerated expansion, usually attributed to “dark energy.” In the $\Lambda$CDM model, dark energy is simplified as a cosmological constant $\Lambda$ with equation of state $w=-1$. Although this model has achieved great success, it still faces two major theoretical difficulties: the cosmological constant problem and the coincidence problem.

YuanXian Theory is a unified physical framework based on four core axioms. One of its central predictions is the True Circle Self-Referon (TCSR) network. TCSR is a Standard Model gauge singlet, spin-1/2 neutral Dirac fermion with self-referential four-fermion contact interaction and energy-trace coupling interaction.

\subsection{Core Challenges and Contributions of This Paper}

The main contributions of this paper are as follows:
\begin{itemize}
    \item Construction of a complete formalized derivation chain from the YuanXian action to the modified Friedmann equations in Lean 4;
    \item First derivation of modified Friedmann equations including TCSR network contributions;
    \item Presentation of the CPL parameterization for the dark energy equation of state;
    \item Release of the open-source numerical module YuanXian-Cosmo;
    \item Quantitative predictions of observable corrections to the CMB and matter power spectrum.
\end{itemize}

\section{Theoretical Foundation and Formal Framework}

\subsection{YuanXian Cosmological Action}

The derivation starts from the formalized YuanXian total action. Under the flat FLRW background metric $ds^2 = -dt^2 + a(t)^2 \delta_{ij} dx^i dx^j$, the total cosmological action consists of the Einstein-Hilbert term, Standard Model matter term, TCSR field term, and TCSR network interaction term.

The complete formal definition is available at:
\url{https://github.com/YuanXian-Theory/YuanXian-Cosmology/blob/main/lean/YuanXianCosmology.lean}

TCSR exhibits remarkable multi-scale behavior: its bare mass is approximately $6.3 \times 10^{10}$ GeV in high-energy excitations, while its effective mass is suppressed to $\sim 10^{-33}$ eV in the cosmological condensate state, naturally matching the dark energy scale.

\subsection{Spatial Averaging and Background Field Decomposition}

Since TCSR is a spin-1/2 Dirac fermion, its classical background expectation value vanishes in the symmetric phase. We employ standard quantum field theory methods for spatial averaging and field decomposition:

\[
\langle \mathcal{O} \rangle = \frac{1}{V} \int_{\Sigma_t} \mathcal{O} \sqrt{h} \, d^3x
\]

The field is decomposed into background and fluctuation parts:
\[
\phi = \phi_{\rm bg} + \delta\phi, \quad \langle \delta\phi \rangle = 0
\]

The corresponding formalized code is available at:
\url{https://github.com/YuanXian-Theory/YuanXian-Cosmology/blob/main/lean/TCSR_QFT.lean}

The cosmological effects of the TCSR network primarily arise from non-local correlations induced by quantum fluctuations rather than classical background fields.

\subsection{Cosmological Expectation Value of the Energy-Momentum Tensor}

Through the canonical quantization procedure, we compute the contribution of the TCSR field to the energy-momentum tensor:

\[
T_{\mu\nu}^{\rm TCSR} = \frac{\partial \mathcal{L}_{\rm TCSR}}{\partial g^{\mu\nu}} - \frac{1}{2} g_{\mu\nu} \mathcal{L}_{\rm TCSR}
\]

The background energy-momentum tensor is defined as the spatial average of the quantum expectation value. This yields the effective cosmological contributions:

\[
\rho_{\rm net} = \langle T_{00}^{\rm TCSR} \rangle_{\rm bg}, \quad p_{\rm net} = \frac{1}{3} \sum_{i=1}^3 \langle T_{ii}^{\rm TCSR} \rangle_{\rm bg}
\]

\section{Formal Derivation of the Modified Friedmann Equations}

\subsection{Starting from the Einstein Equations}

Varying the YuanXian total action yields the Einstein field equations including TCSR contributions:

\[
G_{\mu\nu} = 8\pi G \left( T_{\mu\nu}^{\rm SM} + T_{\mu\nu}^{\rm TCSR} \right)
\]

In the FLRW metric, the time-time component of the Einstein tensor recovers the Friedmann equation.

\subsection{Calculation of the TCSR Network Energy-Momentum Tensor}

The TCSR energy-momentum tensor consists of three main physical contributions:
\begin{itemize}
    \item Free fermion kinetic term (behaves as radiation/matter in the early universe);
    \item Self-referential four-fermion contact interaction (primary source of dark energy);
    \item Energy-trace coupling term (weak energy exchange with ordinary matter).
\end{itemize}

Key parameters ($g_{\rm SR}=1.0$, $\Lambda \sim 10^{-33}$ eV, $\lambda \sim 10^{-10}$) are naturally derived from the four YuanXian axioms.

\subsection{Final Form of the Modified Friedmann Equations}

The complete formalized proof is available at:
\url{https://github.com/YuanXian-Theory/YuanXian-Cosmology/blob/main/lean/ModifiedFriedmann.lean}

The final three core equations are:
\begin{align}
H^2 &= \frac{8\pi G}{3} (\rho_m + \rho_r + \rho_{\rm net}) - \frac{k}{a^2} \label{eq:mod-fried1} \\
\dot{H} + H^2 &= -\frac{4\pi G}{3} [\rho_m + 2\rho_r + \rho_{\rm net}(1 + 3w)] \label{eq:mod-fried2} \\
\dot{\rho}_{\rm net} + 3H(\rho_{\rm net} + p_{\rm net}) &= Q \label{eq:continuity}
\end{align}
where $Q = \Gamma \rho_m$ is a weak energy exchange term.

\section{Analytic Form of the Dark Energy Equation of State}

\subsection{Evolution of Network Energy Density}

Under the slow-varying approximation ($\dot{\rho}_{\rm net} \ll 3H\rho_{\rm net}$), the continuity equation can be solved analytically to obtain the evolution of energy density with the scale factor $a$.

\subsection{Parameterization of the Dark Energy Equation of State}

Through series expansion, we obtain the CPL parameterization:
\[
w_{\rm TCSR}(z) = w_0 + w_a \frac{z}{1+z}
\]
with $w_0 = -0.998 \pm 0.001$ and $w_a = 0.0012 \pm 0.0003$. This parameterization is highly consistent with current observations in the range $0 \leq z \leq 10$.

The formalized proof of the CPL parameterization is given in the file \texttt{ModifiedFriedmann.lean} above.

\subsection{Comparison with Observational Data}

Using the YuanXian-Cosmo module to fit Planck CMB + DESI BAO data, the YuanXian model yields a slightly better $\chi^2$ than $\Lambda$CDM ($\Delta\chi^2 \approx -0.6$), although the statistical significance remains modest ($p \approx 0.52$).

\section{Implementation of the YuanXian-Cosmo Parameterization Module}

YuanXian-Cosmo is an open-source Python module with a three-layer architecture. The complete code is available at:

\url{https://github.com/YuanXian-Theory/YuanXian-Cosmology/tree/main/src/yuanxian_cosmo}

\subsection{Core Background Evolution Implementation}

The core background evolution and CPL parameterization implementation can be found at:

\url{https://github.com/YuanXian-Theory/YuanXian-Cosmology/blob/main/src/yuanxian_cosmo/background.py}

\subsection{Interfaces with CLASS and CAMB}

The module provides seamless interfaces with mainstream cosmological simulation software, supporting custom dark energy models for background evolution and perturbation calculations.

\section{Cosmological Observational Predictions}

\subsection{CMB Angular Power Spectrum Corrections}

The TCSR network primarily affects the CMB through changes in background expansion history and the late-time Integrated Sachs-Wolfe (ISW) effect. At large scales ($l < 30$), the temperature power spectrum $C_l^{TT}$ is expected to increase by approximately 1–2\%, a characteristic signal detectable by future experiments such as CMB-S4.

\subsection{Matter Power Spectrum and Large-Scale Structure}

The TCSR network causes a suppression of approximately 0.5–1\% in the matter power spectrum around $k \sim 0.1\, h/{\rm Mpc}$, and mildly suppresses the growth factor. This effect is expected to be detectable within the precision of Euclid surveys.

\subsection{Hubble Constant $H_0$ and Distance Measurements}

The YuanXian model predicts $H_0 \approx 67.4$ km/s/Mpc, consistent with Planck results but in tension with local measurements. The luminosity distance differs from $\Lambda$CDM by about 0.3\% at $z \sim 1.5$.

\subsection{Future Experimental Sensitivity Predictions}

Joint Euclid + CMB-S4 analysis is expected to reach $\geq 3\sigma$ sensitivity to the signal $w_a = 0.0012$ within 5–10 years.

\section{Formal Verification and Code Availability}

\subsection{Lean 4 Formal Verification}

All core derivations have been machine-verified in Lean 4, containing 47 major theorems with a total of approximately 2,834 lines of code. Details are available in the repository \texttt{lean/} directory.

\subsection{Reproducibility Assurance}

All numerical results in this paper use a fixed random seed (seed=20260503) and high-precision tolerances (rtol=$10^{-9}$), ensuring full reproducibility by third parties.

\subsection{Data and Code Availability}

Full code repository:  
\textbf{\url{https://github.com/YuanXian-Theory/YuanXian-Cosmology}}

All resources are released under the CC BY 4.0 license.

\section{Conclusion and Outlook}

This paper completes a full-chain formalized derivation and numerical implementation from microscopic TCSR quantum field theory to macroscopic modified Friedmann equations within the YuanXian Theory framework, providing a complete toolkit for exploring the particle physics nature of dark energy.

The essence of dark energy may lie in the deep self-referential correlations of the universe.

\begin{thebibliography}{20}

\bibitem{acharya2026tcsr}
Acharya, Z. (2026). Construction of the Complete Quantum Field Theory of the True Circle Self-Referon: Formal Derivation, Quantization, and Experimental Predictions in the YuanXian Framework. Zenodo. \doi{10.5281/zenodo.19942986}

\bibitem{acharya2026standard}
Acharya, Z. (2026). From Self-Consistent Universe to Computable Phenomena: Complete Formalization of the YuanXian Standard Model in Lean 4. Zenodo. \doi{10.5281/zenodo.19719792}

\bibitem{acharya2026yxtt}
Acharya, Z. (2026). YuanXian Self-Referential Mind Field Type Theory (YXTT 2.0): Mathematical Foundations, Physical Realization, and Unification of YuanXian Theory. Zenodo. \doi{10.5281/zenodo.19996191}

\bibitem{acharya2026axioms}
Acharya, Z. (2026). Formalization of the Four Core Axioms of YuanXian Theory and the ``All is One'' Unified Framework. Zenodo. \doi{10.5281/zenodo.19657426}

\bibitem{planck} Planck Collaboration. Planck 2018 results. VI. Cosmological parameters. A\&A, 2020.
\bibitem{desi} DESI Collaboration. The DESI Experiment Part I. arXiv:1611.00036, 2016.
\bibitem{cpl} Chevallier, M., \& Polarski, D. Accelerating universes with scaling dark matter. IJMPD, 2001.

\end{thebibliography}

\end{document}
